\begin{abstract}
Fisher's paired efficiency is $1/(1-\rho)$ when the two margins share a variance.
The companion empirical paper replaces that formula by the Paired Efficiency Factor $\eta = (1+\kappa)/(1+\kappa-2\sqrt{\kappa}\,\rho)$ when the variance ratio is $\kappa$.
This note records the algebraic consequences of that extension, and one sampling consequence.

The substitution $\tau = \tfrac{1}{2}\log\kappa$ gives the canonical form $\eta = \cosh\tau/(\cosh\tau-\rho)$.
The map $\kappa\mapsto 1/\kappa$ leaves $\eta$ unchanged.
The same formula is $\eta = 1/(1-u)$ with $u = \rho\,\sech\tau$.
That scalar is the pairwise intraclass correlation.
It is also Lin's concordance coefficient when the two means coincide.
An affine reciprocal, a Poisson-kernel factor, a latitude reading, and a latent geometric family are remarks on this one scalar.

Under bivariate normality the plug-in variance of $\hat u$ depends on $\tau$ as well as on $u$.
No transformation of $u$ therefore stabilises $\hat\eta$ exactly.
At equal variances the stabilising transform is Fisher's $z = \arctanh\rho$.
The coordinate $\psi = 2\arctanh\sqrt{u}$ stabilises the maximum-likelihood estimator of a latent geometric family on $u>0$.
It does not stabilise the plug-in estimator computed from the pairs.
\end{abstract}
